\Article{18}
\begin{center}\emph{Theorie des biquadratischen Restes $-1-2i$.}\end{center}
Der Modulus $=m=a+bi$ soll so beschaffen sein, dass $a$ ungerade, $b$ gerade; auch setzen wir voraus, dass derselbe eine Primzahl sei.

Der Decident wird durch folgende Schemata vorgestellt, von deren Identität man sich leicht überzeugt:
\begin{center}\resizebox{\textwidth}{!}{\input{diagrams/native332_schemata.tex}}\end{center}

Der Kürze wegen bezeichnen wir $S(x,x+\alpha,x+\alpha+\beta,x+\beta)$ durch $[x,\alpha,\beta]$ so dass
\[
\begin{aligned}
[x,\alpha,\beta]&=-[x,\beta,\alpha]=[x+\alpha,\beta,-\alpha]\\
&=-[x+\alpha,-\alpha,\beta]\\
&=[x+\alpha+\beta,-\alpha,-\beta]\\
&=-[x+\alpha+\beta,-\beta,-\alpha]
\end{aligned}
\]

Setzt man ferner
\[
\frac{m}{-2-4i}=\frac{-a-2b}{10}+\frac{-b+2a}{10}i=Q
\]
\SourceLeaf{333}

so besteht der Decident aus folgenden acht Theilen
\[
\begin{aligned}
\mathrm{I}&=[0,\tfrac12,-iQ]\\
-2\mathrm{II}&=-2[0,\tfrac12,Q]\\
\mathrm{III}&=+[Q,\tfrac12,-iQ]\\
\mathrm{IV}&=+[-iQ,\tfrac12,Q]\\
-3\mathrm{V}&=-3[Q,\tfrac12,Q]\\
-3\mathrm{VI}&=-3[2Q,\tfrac12,-iQ]\\
+2\mathrm{VII}&=+2[(1-i)Q,\tfrac12,Q]\\
+\mathrm{VIII}&=+[0,\tfrac12,\tfrac12im]
\end{aligned}
\]

Ist $F$ indefinite ein Elementarrest des Modulus $-1-2i$, so hat man
\[
\Sigma[2FQ,\tfrac12,iQ]=[0,-\tfrac12-i,\tfrac12im]
\]

Setzt man also für $F$: $0,1,i,-1,-i$ so hat man
\[
\begin{array}{rll}
0={}&[0,\tfrac12,iQ]&=\mathrm{IX}\\
&+[2Q,\tfrac12,iQ]&=\mathrm{X}\pdot\\
&+[2iQ,\tfrac12,iQ]&=\mathrm{XI}\\
&+[-2Q,\tfrac12,iQ]&=\mathrm{XII}\\
&+[-2iQ,\tfrac12,iQ]&=\mathrm{XIII}\\
&-[0,-\tfrac12-i,\tfrac12im]&=-\mathrm{XIV}
\end{array}
\]

Man setze dies zu dem vorigen Werth des Decident hinzu. Aus dieser Vereinigung fliessen folgende Resultate

(1) Da $(1+2i)Q+\frac12$ eine ganze Zahl ist, so wird
\[
\begin{aligned}
\mathrm{XI}&=[-Q-\tfrac12,\tfrac12,iQ]= -[-Q,-\tfrac12,iQ]\\
&=-[Q,\tfrac12,-iQ]
\end{aligned}
\]
also
\[
\mathrm{III}+\mathrm{XI}=0
\]

(2) Wir ziehen zusammen $\mathrm{IV},-3\mathrm{V},\mathrm{X},\mathrm{XIII}$ auf folgende Weise
\[
\begin{aligned}
\mathrm{IV}&=+[-2Q+\tfrac12i,\tfrac12,Q]=[-2iQ-\tfrac12,\tfrac12i,iQ]\\
\mathrm{V}'&=-[2Q,+\tfrac12,-Q]=-[-2iQ,-\tfrac12i,iQ]\\
\mathrm{X}&=[+iQ-\tfrac12i,\tfrac12,iQ]\\
&=[-iQ+\tfrac12i,-\tfrac12,-iQ]=[-2iQ-\tfrac12+\tfrac12i,\tfrac12,iQ]\\
\mathrm{XIII}&=[-2iQ,\tfrac12,iQ]
\end{aligned}
\]

\SourceLeaf{334}

Also die ganze Ausbeute aus diesen Theilen
\[
\begin{aligned}
&-4\mathrm{V}\\
&+R(-2iQ)-R(-iQ)-I(-2iQ)+I(-iQ)\\
&-\text{Quadr. }[-2iQ-\tfrac14+\tfrac14i]\\
&+\text{Quadr. }[-2iQ+\tfrac14-\tfrac14i]
\end{aligned}
\]

(3) $\mathrm{I},-2\mathrm{II},-3\mathrm{VI},+2\mathrm{VII},+\mathrm{IX},+\mathrm{XII}$ zusammengezogen geben folgendes
\[
\begin{aligned}
\mathrm{I}&=[0,\tfrac12,-iQ]\\
\mathrm{II}&=[0,-\tfrac12i,-iQ]\\
\mathrm{VI}&=[iQ-\tfrac12i,\tfrac12,-iQ]=[\tfrac12-\tfrac12i,-\tfrac12,iQ]\\
&=[-\tfrac12+\tfrac12i,\tfrac12,-iQ]\\
\mathrm{VII}&=[-Q-\tfrac12i,\tfrac12,Q]=[-\tfrac12+\tfrac12i,-\tfrac12,-Q]\\
&=[+\tfrac12+\tfrac12i,-\tfrac12i,-iQ]\\
\mathrm{IX}&=[0,-\tfrac12,-iQ]\\
\mathrm{XII}&=[-iQ+\tfrac12i,\tfrac12,iQ]=[\tfrac12+\tfrac12i,-\tfrac12,-iQ]
\end{aligned}
\]
\[
\begin{aligned}
={}&+4\mathrm{I}-4\mathrm{II}-4\mathrm{VI}+4\mathrm{XII}\\
&-I(\tfrac12i)+I(-iQ+\tfrac12i)+I0-I(-iQ)-2R0+2R(-iQ)\\
&+2\text{Quadr. }(-iQ+\tfrac14+\tfrac14i)\\
={}&4\mathrm{I}-4\mathrm{II}-4\mathrm{VI}+4\mathrm{XII}\\
&+R(-1-2i)Q-R(-2iQ)-I(-iQ)+2R(-iQ)\\
&+2\text{Quadr. }(-2iQ+\tfrac14+\tfrac14i)
\end{aligned}
\]

Dies Alles zusammen gibt folglich
\[
\begin{aligned}
&+4\mathrm{I}-4\mathrm{II}-4\mathrm{V}-4\mathrm{VI}+4\mathrm{XII}+\tfrac12(a-1)-\tfrac12b\\
&+\text{Quadr. }(-2iQ+\tfrac14-\tfrac14i)\\
&+2\text{Quadr. }(-2iQ+\tfrac14+\tfrac14i)\\
&-\text{Quadr. }(-2iQ-\tfrac14+\tfrac14i)
\end{aligned}
\]

Endlich gibt $\mathrm{VIII}-\mathrm{XIV}=\frac12(a-1)+\frac12b$
Also da die drei Quadrattheile dem Decident von $\frac{m}{-1-2i}$ gleich sind, so wird
\[
\text{Dec. }\frac{-1-2i}{m}\equiv a-1+\text{Dec. }\frac{m}{-1-2i}\quad\text{W. Z. B. W.}
\]

Wahrscheinlich wird der Beweis noch sehr dadurch vereinfacht werden können, dass
\[
\text{Dec. }\frac{1+2i}{m}=\vcenter{\hbox{\resizebox{.35\textwidth}{!}{\input{diagrams/native334_triangle.tex}}}}
\]

\SourceLeaf{335}

\Article{19}
Durch Induction ist folgendes gefunden
\[
\begin{aligned}
\text{Dec. }\frac{a-bi}{a+bi}&\equiv\frac{aa+2ab-1}{4},& a&\equiv1\ (\mathrm{mod.}\ 4)\\
&\equiv\frac{aa+2ab+2bb-1}{4},&a+bi&\equiv1\ (\mathrm{mod.}\ 2+2i)
\end{aligned}
\]

Hiemit steht Folgendes in Verbindung:
Es sei $aa+bb=p$ (Primzahl) $a\equiv1\ (\mathrm{mod.}\ 4)$
\[
\begin{aligned}
1\pdot2\pdot3\mathinner{.\,.\,.\,.\,.\,.\,.\,.\,.\,.\,.}\tfrac14(p-1)&\equiv\alpha\ \mathrm{mod.}p\\
\tfrac14(p+3)\pdot\tfrac14(p+7)\ldots\tfrac12(p-1)&\equiv\beta\\
\tfrac12(p+1)\pdot\tfrac12(p+3)\ldots\tfrac34(p-1)&\equiv\gamma\\
\tfrac14(3p+1)\mathinner{.\,.\,.\,.\,.\,.\,.\,.\,.}\ p-1&\equiv\delta
\end{aligned}
\]
so ist
\[
\begin{aligned}
\alpha&\equiv\delta,&\beta&\equiv\gamma&\text{wenn }\tfrac12b\text{ gerade}\\
\alpha&\equiv-\delta,&\beta&\equiv-\gamma&\text{wenn }\tfrac12b\text{ ungerade}
\end{aligned}
\]
\[
\begin{gathered}
\pm\alpha\beta\equiv i,\quad\pm\beta\beta\equiv2b,\quad\frac\beta\alpha\equiv2a,\\
\frac\beta{1+\alpha\beta}\equiv\surd a,\quad\tfrac12\beta(1-\alpha\beta)\equiv\surd a
\end{gathered}
\]

Es wird demnach nur darauf ankommen die Decidenten bei reellen Resten zu bestimmen
\[
\begin{gathered}
a^{\frac14bb}\pdot b^{\frac14(aa-1)}\equiv1\ (\mathrm{mod.}\ (aa+bb))\\
\text{si }a\equiv1)(\mathrm{mod.}\ 4)\quad b\text{ par}\quad aa+bb\text{ primus}
\end{gathered}
\]

Will man bloss mit reellen Zahlen zu thun haben, so kommt es auf folgendes Haupttheorem an. Es sei $a-1$ durch $4$, $b$ durch $2$ theilbar; $a$ und $b$ ohne gemeinschaftlichen Divisor, $k$ bedeute die Zahlen $1,2,3\ldots aa+bb-1$.
Es sei
\begin{list}{}{\setlength{\labelwidth}{1.5em}\setlength{\leftmargin}{2em}}
\item[$\alpha$] die Zahl aller Werthe von $k$, wo die kleinsten Reste von $ak,\allowbreak bk,\quad\allowbreak aak,\quad\allowbreak abk$ alle zwischen $0$ und $\frac12(aa+bb)$ liegen
\item[$\beta$] $ak,bk,\quad aak,-abk$
\item[$\gamma$] $ak,bk,-aak,-abk$
\item[$\delta$] $ak,bk,-aak,\quad abk$
\end{list}
alsdann ist $\beta+2\gamma+3\delta-\frac14(aa-1)$ durch $4$ theilbar.
\BookRule{.35\textwidth}
